% -----------------------------------------------------------------------------------
\section{The witness problems}\label{sec:witnesses}

This section works through the problems that separate the criteria. Payoffs are scaled for
arithmetic convenience; nothing depends on the particular numbers beyond the inequalities stated.

\subsection{Newcomb's problem}\label{sec:newcomb}

Let $O$ denote taking only the opaque box and $B$ taking both boxes. A predictor of symmetric error
$\varepsilon\in[0,1]$ forecasts the agent's prescription at its single information set, which is
reached with probability one; the opaque box is full if and only if the forecast is $O$. Payoffs are

\begin{center}
\small
\begin{tabular}{@{}lcc@{}}
  \toprule
  & opaque box full & opaque box empty \\
  \midrule
  $O$ & 2 & 0 \\
  $B$ & 3 & 1 \\
  \bottomrule
\end{tabular}
\end{center}
with $\Pr(\text{full}\mid O)=1-\varepsilon$ and $\Pr(\text{full}\mid B)=\varepsilon$. The structural
states are the two contents of the opaque box.

\begin{proposition}[\cite{nozick1969}]\label{prop:newcomb}
In Newcomb's problem $V(O)=2(1-\varepsilon)$ and $V(B)=1+2\varepsilon$, so $V(O)-V(B)=1-4\varepsilon$
and \cri{1} uniquely requires $O$ whenever $0\le\varepsilon<\tfrac14$. Since $U(s,B)-U(s,O)=1$ in both
structural states, \cri{2} and \cri{3} uniquely require $B$ for every $\varepsilon$. Hence no theory
satisfies \cri{1} and \cri{2} on any class containing Newcomb's problem with $\varepsilon<\tfrac14$.
\end{proposition}

With opaque prize $m$ and transparent bonus $c$, the same conflict holds exactly when
$m(1-2\varepsilon)>c$. Any symmetric predictor better than chance therefore produces it for suitable
stakes. The conflict is Nozick's original one between maximising expected utility and dominance;
the content of Proposition~\ref{prop:newcomb} is only that the formal criteria reproduce it.

The proposition shows that a conflict \emph{can} arise. The next result characterises exactly when
it can, and shows that neither accurate prediction nor large stakes is needed.

\begin{theorem}[Policy-dependent state laws characterise the conflict]\label{thm:characterisation}
Let $A$ be a finite set of actions with $|A|\ge2$, $S$ a finite set of states, and $\mu_a\in\Delta(S)$
the state law induced by prescribing $a$ at a single always-reached choice. The following are
equivalent:
\begin{enumerate}[label=(\roman*)]
  \item for every payoff function $u\colon A\times S\to[0,1]$, some ex-ante maximising action is not
        strictly dominated state by state;
  \item the laws $\mu_a$, $a\in A$, are all equal.
\end{enumerate}
\end{theorem}

\begin{proof}
(ii)$\Rightarrow$(i). Under a common law $\mu$, if $a$ strictly dominates $b$ in every state then
$\E_\mu u(a,\cdot)>\E_\mu u(b,\cdot)$, so no ex-ante maximiser is strictly dominated; a maximiser
exists because $A$ is finite.

(i)$\Rightarrow$(ii). Suppose $\mu_a\neq\mu_b$. There is an event $E\subseteq S$ on which they differ;
exchanging the names of $a$ and $b$ if necessary,
$\delta=\mu_b(E)-\mu_a(E)>0$. Choose $0<c<\delta/(1+\delta)$ and set
\[
  u(b,s)=(1-c)\,\one_E(s),\qquad u(a,s)=(1-c)\,\one_E(s)+c,
\]
with payoff $0$ for every other action. All payoffs lie in $[0,1]$, and $a$ strictly dominates $b$ by
$c$ in every state. Yet
\[
  V(b)-V(a)=(1-c)\mu_b(E)-(1-c)\mu_a(E)-c=(1-c)\delta-c>0,
\]
and $V(a)\ge c>0$ exceeds the value $0$ of every other action. So $b$ is the unique ex-ante
maximiser, and it is strictly dominated.
\end{proof}

Any non-constant dependence of the preceding state law on the prescription thus suffices for some
bounded payoff assignment to separate \cri{1} from \cri{2}; for fixed payoffs, the size of the
dependence determines whether the conflict occurs.

\subsection{The twin prisoner's dilemma}\label{sec:twin}

Two copies of the same policy play a prisoner's dilemma with temptation $3$, reward $2$, punishment
$1$, and sucker's payoff $0$. The twin is represented as a prediction node that copies the agent's
prescription at its single information set; it is \emph{not} a second, independently assignable
coordinate of $\Pol_D$, which would remove the stipulated identity of the two decisions. Then
$V(\text{cooperate})=2$ and $V(\text{defect})=1$, so \cri{1} requires cooperation, while with the
twin's action held fixed, defection gains one unit whether the twin cooperates or defects, so
\cri{2} and \cri{3} require defection. The connection between this problem and Newcomb's is due to
\textcite{lewis1979}. What the tree does not supply is the identity of the two decisions:
\enquote{the same computation} enters only as a stipulated dependence of the prediction node on
the agent's prescription. Sections~\ref{sec:inference} and~\ref{sec:operators} ask whether an agent
given the twin's source code could infer that dependence, and Section~\ref{sec:counterpossibles}
records how each candidate answer decides this problem.

\subsection{Death in Damascus}\label{sec:did}

Death predicts which of two cities the agent will choose and waits there
\autocite{gibbardharper1978}. Let actions and states both be $\{1,2\}$, with
$U(s,a)=\one\{a\neq s\}$ and $\mu_a=\delta_a$ (infallible prediction).

\begin{proposition}\label{prop:did}
In Death in Damascus no deterministic policy satisfies \cri{3}, whereas both deterministic policies
satisfy \cri{2}. The failure of \cri{3} persists for every symmetric error $\varepsilon<\tfrac12$.
\end{proposition}

\begin{proof}
Under the policy $a$, the induced state is $a$; staying has causal value $0$ and switching has causal
value $1$. With error $\varepsilon$, staying has causal value $\varepsilon$ and switching $1-\varepsilon$,
and $\varepsilon<1-\varepsilon$. For \cri{2}, neither city is better than the other in \emph{both}
structural states, so no action is dominated.
\end{proof}

Consequently \cri{3}, read as a requirement on deterministic policies, is individually
unsatisfiable on any class containing this problem: it fails to define a decision theory at all.
The absence of ratifiable pure acts in this problem is well known from the ratifiability
literature \autocite{skyrms1990ratif}; Section~\ref{sec:mixedexist} records how far randomisation
repairs it.

\subsection{The Bomb}\label{sec:bomb}

\textcite{macaskill2019} describes a predictor of error $10^{-24}$ that places a bomb in the left
box if and only if it predicts that the agent will take the right box, and leaves a note saying so.
The agent sees the note and the bomb. Taking the left box is free; taking the right box costs $1$;
burning to death costs $M=10^{9}$ in the scaling used here.

Once the agent observes whether a bomb is present, there are \emph{two} information sets,
$h_{\mathrm{bomb}}$ and $h_{\mathrm{clear}}$, and four policies
$(a_{\mathrm{bomb}},a_{\mathrm{clear}})\in\{L,R\}^2$. This has a consequence for well-definedness.

\begin{proposition}[The Bomb as stated is not well defined]\label{prop:bombill}
No predictor whose forecast determines the observation (\enquote{predicts $R$: bomb; predicts $L$:
no bomb}) and whose target is the agent's \emph{eventual} action can have an error bound below one
uniformly over the four policies.
\end{proposition}

\begin{proof}
Under the policy $(L,R)$, a forecast of $R$ produces a bomb and hence the action $L$; a forecast of $L$
produces no bomb and hence the action $R$. The forecast is wrong in every case.
\end{proof}

A coherent version is obtained by letting the predictor target the fixed coordinate
$\pi(h_{\mathrm{bomb}})$, whether or not $h_{\mathrm{bomb}}$ is eventually reached:
\[
  \Pr\bigl(\text{bomb}\mid a_{\mathrm{bomb}}=L\bigr)=\varepsilon,\qquad
  \Pr\bigl(\text{bomb}\mid a_{\mathrm{bomb}}=R\bigr)=1-\varepsilon.
\]

\begin{proposition}\label{prop:bomb}
In the repaired Bomb with death cost $M>1$, the ex-ante values are
\begin{center}
\small
\begin{tabular}{@{}lc@{}}
  \toprule
  policy $(a_{\mathrm{bomb}},a_{\mathrm{clear}})$ & $V$ \\
  \midrule
  $(L,L)$ & $-M\varepsilon$ \\
  $(L,R)$ & $-M\varepsilon-(1-\varepsilon)$ \\
  $(R,L)$ & $-(1-\varepsilon)$ \\
  $(R,R)$ & $-1$ \\
  \bottomrule
\end{tabular}
\end{center}
For $0<\varepsilon<1/(M+1)$, \cri{1} uniquely selects $(L,L)$, while \cri{2} and \cri{3} uniquely
select $(R,L)$.
\end{proposition}

\begin{proof}
The values follow by conditioning on the presence of the bomb. $(L,L)$ beats $(R,L)$ exactly when
$M\varepsilon<1-\varepsilon$, that is $\varepsilon<1/(M+1)$; it beats $(L,R)$ and $(R,R)$ a fortiori. For
$\varepsilon>0$ both information sets are positively reached under every policy. At
$h_{\mathrm{bomb}}$, every structural state reaching it contains the bomb, and $R$ yields $-1>-M$;
at $h_{\mathrm{clear}}$, $L$ yields $0>-1$. Hence $(R,L)$ is the only policy satisfying \cri{2},
and, since it is causally optimal at both information sets, it satisfies \cri{3}.
\end{proof}

MacAskill's numbers satisfy the inequality. Two qualifications are essential. At $\varepsilon=0$,
$h_{\mathrm{bomb}}$ is unreached under $(L,L)$, and the reached-only criteria \cri{2} and \cri{3} do
not oppose that policy. And if the problem is instead taken to be the already-conditioned decision
\enquote{a bomb is visibly present}, its two options have values $-M$ and $-1$, and \cri{1} itself
selects $R$. The visible-bomb challenge is MacAskill's; the complete contingent-policy kernel above
is the additional assumption under which it becomes a theorem.

\subsection{Counterfactual mugging}\label{sec:mugging}

A fair coin is tossed. On heads the agent receives $r$ if and only if an infallible predictor
forecasts that the agent's prescription on tails is to pay; on tails the agent is asked to pay $c$,
with $r>c>0$ \autocite{nesov2009}. There is one information set, reached on tails. Then
\[
  V(\text{Pay})=\tfrac12(r-c)>0=V(\text{Refuse}),
\]
so \cri{1} requires paying, while at the tails information set paying costs $c$ in every structural
state reaching it, so \cri{2} and \cri{3} require refusing. The problem separates ex-ante from
ex-post evaluation by a prediction about an information set that is unreached on half the paths; it
is used in Section~\ref{sec:evidential} to locate evidential criteria and in
Section~\ref{sec:targets} to test restrictions on what predictors may target.
